\section{Suma del error de los menores reales sin contar tipos}
Sean $\sqrt n\le j\le n$, $5\le v\le8\log n$ y $q=4v+12$. En cada soporte $U$
de $v$ casillas, se retiran de la máscara basal todas las casillas de $U$ para definir
un exterior común, también en las caras que desactivan algunas restricciones diagonales.
Cada colocación de torres que cuenta en la cara $Y\subseteq U$ tiene un único conjunto
$I$ de sitios ocupados en $U$. Se permiten los $I$ compatibles con RC y con los ataques
activados en $Y$. Su peso es el permanente del exterior con las filas y columnas de $I$
suprimidas, por los factores exactos de sitio y escalado si se usan matrices ponderadas.
No se imponen ataques diagonales externos en este experimento de activación parcial.
El átomo vacío es positivo cuando el exterior tiene emparejamiento, como en la
clase de menores uniformes. Los átomos imposibles tienen peso cero en ambos lados.
Los menores positivos utilizados han de admitir el escalado controlado de la parte 02
de esa fuente. La existencia de un único emparejamiento no implica por sí sola ese escalado.

\paragraph{Lema (condicionado a la entrada real de menores).}
Supongamos que cada peso físico positivo $W_I$ se sustituye por un peso positivo
$\widehat W_I$ con los mismos factores de escalado y
\[
 |\log(\widehat W_I/W_I)|\le \eta_v,\qquad
 \eta_v=(C'q/j)^q+C'e^{-cj}.
\]
Es la precisión de \emph{Menores uniformes}, partes 01, 04 y 05, con $h=|I|\le v$
y $j-h\ge j/2$; una normalización aproximada común se puede absorber en las constantes.
El orden requerido llega a $4\lfloor8\log n\rfloor+12$; cabe en el umbral lineal
de esa fuente para $j\ge\sqrt n$ y $n$ suficientemente grande.
Sea $F(Y)$ el logaritmo del conteo de la cara dividido por el permanente basal.
Sea $\widehat F_{U,q}(Y)$ la misma expresión con esos átomos. El denominador exacto
se conserva, y toda constante común a las caras desaparece en la proyección $\Delta_U$.
En cada cara, al sumar $e^{-\eta_v}W_I\le\widehat W_I\le e^{\eta_v}W_I$,
\[
 |\widehat F_{U,q}(Y)-F(Y)|\le\eta_v,\qquad
 |\Delta_U(\widehat F_{U,q}-F)|\le 2^v\eta_v .
\]
Aquí $\Delta_U f=\sum_{Y\subseteq U}(-1)^{|U|-|Y|}f(Y)$. Cada soporte tiene
su propio aproximante; la estimación no necesita una función global aproximada.
Se cuentan a lo sumo $j^{2v}$ soportes, sin clasificación geométrica. Para $n$ grande
se tiene $C'q\le j^{1/3}$ simultáneamente en el intervalo de grados. Luego
\[
 2^v j^{2v}(C'q/j)^q\le 2^v j^{-8-2v/3}.
\]
La suma de la derecha es $O(j^{-8})$. Para la parte exponencial el número total,
incluidas las caras, es $\exp(O((\log n)^2))$. Al multiplicar por $e^{-cj}$ se
obtiene $O(j^{-8})$ también. En consecuencia,
\[
 \sum_{5\le |U|\le8\log n}|\Delta_U(\widehat F_{U,q}-F)|=O(j^{-8}),
 \qquad q=4|U|+12 .
\]
El resultado vale en cada máscara con la entrada uniforme. Sirve también en cualquier
subprobabilidad que sólo sume tales máscaras. No se aplica al error de la parte
canónica $\Lambda$, al propio aproximante, ni a grados superiores a $8\log n$.
